\documentclass[11pt]{article}
\usepackage[T1]{fontenc}
\usepackage{lmodern}
\usepackage[margin=1in]{geometry}
\usepackage{amsmath,amssymb,amsthm,booktabs,microtype}
\usepackage[hidelinks]{hyperref}
\newtheorem{theorem}{Theorem}
\newtheorem{lemma}[theorem]{Lemma}
\newtheorem{proposition}[theorem]{Proposition}
\theoremstyle{remark}
\newtheorem{remark}[theorem]{Remark}
\DeclareMathOperator{\Pf}{Pf}
\DeclareMathOperator{\sgn}{sgn}
\newcommand{\E}{\mathbb E}
\newcommand{\Prb}{\mathbb P}
\newcommand{\R}{\mathbb R}
\newcommand{\Z}{\mathbb Z}
\title{Diagonal asymptotics of square Hardinian arrays\\[4pt]
\large A formula for every fixed parameter}
\author{Aurantium}
\date{14 October 2026}
\begin{document}
\maketitle
\begin{abstract}
Starting from the complementary-minor identity of Dougherty-Bliss and Kauers
and the sum-of-squares reduction in the public $r=2,3$ note, we prove
\[
 H_r(n,n)\sim c_r\,4^{rn}n^{-\binom r2},\qquad
 c_r=\frac{2^{r(r-3)}}{3^{r^2}\pi^r}
       \prod_{j=1}^r\Gamma(j/2)^2,
\]
for every fixed integer $r\ge1$. In particular,
$H_4(n,n)\sim4\cdot256^n/(3^{16}\pi^2 n^6)$, in agreement with the published asymptotic table.
The proof uses a Pfaffian sign identity, a finite-order local binomial expansion,
alternating-polynomial cancellation, a Gaussian integral, and an explicit
geometric tail bound. We also give a direct calculation of the four-row constant
which does not require the Gaussian integral evaluation. The formula matches
the published constants through $r=5$, but its $r=7$ specialization disproves
the separate suggestion that every constant is a product of integer powers
of $2,3,\pi$. The results concern $r$ fixed before $n$ tends to infinity.
\end{abstract}

\section{Problem, scope, and statement}

For integers $r\ge0$ and $n\ge1$, let $H_r(n,n)$ count arrays
$(a_{ij})_{1\le i,j\le n}$ with entries in $\Z_{\ge0}$ such that
\[
 a_{11}=0,\qquad a_{nn}=n-r-1,
 \qquad |a_{ij}-(\max(i,j)-1)|\le r,
\]
and each permitted step right, down, or diagonally down-right increases the
entry by $0$ or $1$. These are the square arrays defined in~\cite{dk}.
The asymptotic table on page~12 of its 13-page journal version gives the
fixed-parameter cases through $r=5$. The source identity used below occurs
in the proof of its Theorem~6, on page~9. This note extends the public
$r=2,3$ treatment~\cite{prior} to every fixed parameter.

For $n\le r$, the prescribed endpoint is negative and the count is zero. For
$n=r+1$ the all-zero array is the unique array. In particular, $H_4(n,n)=0$ for
$n=1,2,3,4$ and $H_4(5,5)=1$. These boundary cases do not change an asymptotic
statement. All limits here are through integer $n$ with $r$ fixed first. The
proof uses classical finite combinatorics and real analysis.

\begin{theorem}\label{thm:main}
For each fixed integer $r\ge1$, put $R=\binom r2$. Then
\begin{equation}\label{eq:main}
 H_r(n,n)\sim
 \frac{2^{r(r-3)}}{3^{r^2}\pi^r}
 \left(\prod_{j=1}^r\Gamma(j/2)^2\right)4^{rn}n^{-R}.
\end{equation}
For $r=0$, $H_0(n,n)=1$ exactly. In particular,
\begin{equation}\label{eq:target}
 \boxed{H_4(n,n)\sim\frac{4\,256^n}{3^{16}\pi^2n^6}.}
\end{equation}
The same theorem also proves the $r=5$ row of the source table:
\[
 H_5(n,n)\sim\frac{16\,1024^n}{3^{23}\pi^2n^{10}}.
\]
\end{theorem}

The sum-of-squares identity and the probability interpretation in the next
section are established starting material from the public note~\cite{prior}.
They are reproduced to fix all conventions and dependencies. No novelty is
claimed for those identities, for Pfaffian algebra, for Stirling's expansion,
or for the Gaussian Selberg integral.

\section{Exact finite identities}

Fix $n\ge r+1$ and put $N=n-2$. For a binomial coefficient outside its usual
integer range, use the value zero. Rows and columns in every submatrix are in
increasing order, and an empty determinant is $1$. Write
\[
 M=\left[\binom{i+j}{i}\right]_{0\le i,j\le N}.
\]
The counting identity from~\cite{dk} is
\begin{equation}\label{eq:source}
 H_r(n,n)=\sum_{\substack{I,J\subseteq\{0,\ldots,N\}\\|I|=|J|=r}}
              \det M_{I^c,J^c}.
\end{equation}
Every subset occurs once; there is no factorial multiplicity.

\begin{lemma}[Sum of squares and uniform bound]\label{lem:squares}
For $0\le u_1<\cdots<u_r\le N$, define
\begin{align}
 W_r(u_1,\ldots,u_r)
 &=\sum_{0\le i_1<\cdots<i_r\le N}
       \det\left[\binom{u_a}{i_b}\right]_{a,b=1}^r,\label{eq:W}\\
 w_r(u_1,\ldots,u_r)&=2^{-\sum_a u_a}W_r(u_1,\ldots,u_r).
\end{align}
Then
\begin{equation}\label{eq:squares}
 H_r(n,n)=\sum_{0\le u_1<\cdots<u_r\le N}W_r(u_1,\ldots,u_r)^2.
\end{equation}
If the $X_a$ are independent with $X_a\sim\operatorname{Bin}(u_a,1/2)$, then
\begin{equation}\label{eq:signexpectation}
 w_r(u_1,\ldots,u_r)=
 \E\left[\sgn\prod_{a<b}(X_b-X_a)\right],\qquad |w_r|\le1.
\end{equation}
Here $\sgn(0)=0$.
\end{lemma}
\begin{proof}
Set $P_{u,i}=\binom ui$ and $E=\operatorname{diag}((-1)^i)$. Vandermonde's
binomial identity gives $M=PP^T$. Since $P$ is unit lower triangular and
$P^{-1}=EPE$, we have $\det M=1$ and
\[
 M^{-1}=E(P^TP)E.
\]
Jacobi's complementary-minor identity therefore gives
\[
 \det M_{I^c,J^c}
 =(-1)^{\sum I+\sum J}\det(M^{-1})_{J,I}
 =\det(P^TP)_{J,I}.
\]
All matrices are finite and $M$ is invertible, as required. Cauchy--Binet
expands the last determinant as a sum over $r$-element row sets of $P$.
Summing over $I,J$ gives~\eqref{eq:squares}. Expanding~\eqref{eq:W} and sorting
the distinct values of the independent $X_a$ gives~\eqref{eq:signexpectation}.
A tie contributes zero to the sign and has no corresponding strictly increasing
column set. Since the sign has absolute value at most one, the bound follows.
\end{proof}

For $u<v$, define
\begin{equation}\label{eq:Q}
 D(u,v)=\sum_{j=u}^{v-1}\binom{u+v}{j},\qquad
 Q(u,v)=2^{-u-v}D(u,v),
\end{equation}
and extend $D,Q$ skew-symmetrically, with diagonal entries zero. For independent
$X\sim\operatorname{Bin}(u,1/2)$ and $Y\sim\operatorname{Bin}(v,1/2)$,
$Y-X$ has the distribution of $Z-u$ for $Z\sim\operatorname{Bin}(u+v,1/2)$.
Symmetry about $(u+v)/2$ gives
\[
 \E\sgn(Y-X)=\Prb(Z>u)-\Prb(Z<u)=\Prb(u\le Z<v)=Q(u,v).
\]

\begin{lemma}[Exact Pfaffian reduction]\label{lem:pf}
Let $A_{ij}=Q(u_i,u_j)$ for $1\le i,j\le r$. If $r$ is even, then
$w_r=\Pf A$. If $r$ is odd, then
\begin{equation}\label{eq:augmented}
 w_r=\Pf\begin{pmatrix}A&\mathbf 1\\-\mathbf 1^T&0\end{pmatrix}.
\end{equation}
In particular,
\begin{equation}\label{eq:pf4}
 w_4=Q(u_1,u_2)Q(u_3,u_4)-Q(u_1,u_3)Q(u_2,u_4)
                         +Q(u_1,u_4)Q(u_2,u_3).
\end{equation}
\end{lemma}
\begin{proof}
For any even number of real numbers $x_1,\ldots,x_{2m}$,
\begin{equation}\label{eq:signpf}
 \Pf[\sgn(x_j-x_i)]_{i,j=1}^{2m}
   =\sgn\prod_{i<j}(x_j-x_i).
\end{equation}
If two values coincide, the corresponding rows and columns coincide and both
sides are zero. If all values are distinct, permuting them into increasing order
multiplies both sides by the same permutation sign. In increasing order the
matrix has every entry above the diagonal equal to one. Its Pfaffian is one,
by induction using the first-row expansion and
$\sum_{j=2}^{2m}(-1)^j=1$. This proves~\eqref{eq:signpf}, including ties.
For odd $r$, append a formal value $+\infty$, whose pairwise signs with the
original variables are all one. In a Pfaffian matching each random variable
appears in exactly one factor. Independence therefore factors the expectation
of each product into pair expectations. Apply~\eqref{eq:signexpectation} and
\eqref{eq:Q}; the dummy factor in the odd case is deterministically one.
\end{proof}

\section{A uniform local expansion to arbitrary fixed order}

We need coefficient structure, not a central limit theorem without an error
estimate. All constants in this section may depend on the fixed expansion order
$K$ and the fixed number $C$, but not on $N,L,a,b$ in the specified ranges.

\begin{lemma}\label{lem:local}
Fix $K\ge0$ and $C>0$. There are polynomials $P_k(a,b)\in\mathbb Q[a,b]$,
$0\le k\le K$, such that, uniformly for integer offsets
$0\le b<a\le L\le C\log N$,
\begin{equation}\label{eq:local}
 Q(N-a,N-b)=\frac1{\sqrt{\pi N}}\sum_{k=0}^K\frac{P_k(a,b)}{N^k}
 +O_{K,C}\left(\frac{(1+L)^{2K+3}}{N^{K+3/2}}\right).
\end{equation}
Every $P_k$ is antisymmetric in $a,b$, has total degree at most $2k+1$, and its
homogeneous part of degree $2k+1$ is
\begin{equation}\label{eq:top}
 [P_k(a,b)]_{2k+1}
 =\frac{(-1)^k(a-b)^{2k+1}}{4^k k!(2k+1)}.
\end{equation}
In particular $P_0(a,b)=a-b$.
\end{lemma}
\begin{proof}
Put $d=a-b$, $s=a+b$, and $m=2N-s$. The summation indices in~\eqref{eq:Q}
are $j=m/2+t$, with
\[
 t=-d/2,-d/2+1,\ldots,d/2-1.
\]
The values $t$ may be half-integers; the factorial arguments are always integers.
For sufficiently large $N$, all three factorial arguments $m,m/2+t,m/2-t$
are positive and comparable to $N$, uniformly in the window.

Here is the finite-order Stirling calculation including the degree information.
Write $x=2t/m$ and $c_h=B_{2h}/(2h(2h-1))$. For an integer $J$ with
$2J+1\ge K+1$, the logarithmic Stirling formula with its positive-real-axis
remainder bound~\cite{dlmf} gives
\begin{align}
 &\log\left(2^{-m}\binom{m}{m/2+t}\sqrt{\pi N}\right)\notag\\
 &=-\frac12\log\left(1-\frac{s}{2N}\right)
   -\frac m2\bigl((1+x)\log(1+x)+(1-x)\log(1-x)\bigr)
   -\frac12\log(1-x^2)\notag\\
 &\quad+\sum_{h=1}^J c_h
 \left(m^{-(2h-1)}-(m/2+t)^{-(2h-1)}-(m/2-t)^{-(2h-1)}\right)
 +O_J(N^{-2J-1}).\label{eq:logmass}
\end{align}
One may increase $J$ freely before this finite expansion. All arguments of
Stirling's formula are positive reals, so its standard remainder bound applies;
there is no large complex shift or unverified sector condition.

Taylor-expand~\eqref{eq:logmass} in $N^{-1}$ through order $K$.
The coefficient $\ell_j(s,t)$ at $N^{-j}$ is an even polynomial in $t$ of
total degree at most $j+1$, and
\[
 \ell_1(s,t)=-t^2+s/4-1/8.
\]
To verify the degree bound explicitly, an entropy term is a multiple of
$t^{2h}m^{-(2h-1)}$. Its $N^{-j}$ coefficient is a multiple of
$t^{2h}s^{j-2h+1}$, of degree $j+1$. The logarithmic prefactors give degree
at most $j$, and the $h$th Stirling correction gives degree at most $j-2h+1$.
The symmetry in $t$ follows directly from~\eqref{eq:logmass}.

The remainders of these Taylor expansions are uniform because
$|s|/N,|t|/N\to0$ uniformly and $s,|t|=O(\log N)$. More precisely they are
$O_{K,C}((1+L)^{K+2}N^{-K-1})$ before exponentiation. Expanding the exponential
now yields
\begin{equation}\label{eq:massgeneral}
 2^{-m}\binom m{m/2+t}=
 \frac1{\sqrt{\pi N}}\left(
  \sum_{k=0}^K\frac{T_k(s,t)}{N^k}
  +O_{K,C}\left(\frac{(1+L)^{2K+2}}{N^{K+1}}\right)\right).
\end{equation}
Indeed, the whole logarithm is $O((1+L)^2/N)$, and a Taylor remainder of
order $K+1$ in its exponential has precisely the displayed bound.
The polynomial $T_k$ is even in $t$ and has total degree at most $2k$.
Its degree-$2k$ part is uniquely $(-t^2)^k/k!$: an $\ell_j$ with $j>1$
has degree $j+1<2j$, and the other terms of $\ell_1$ have degree below two.

For each $h\ge0$, Faulhaber's identity gives
\begin{equation}\label{eq:Faulhaber}
 \sum_{v=0}^{d-1}(v-d/2)^{2h}
 =\frac{B_{2h+1}(d/2)-B_{2h+1}(-d/2)}{2h+1}.
\end{equation}
This is an odd polynomial in $d$ of degree at most $2h+1$, with leading term
$d^{2h+1}/(4^h(2h+1))$. Summing each monomial of $T_k(s,t)$ over this grid
therefore gives an odd polynomial in $d$, with $s$ unchanged, of total degree
at most $2k+1$. It is antisymmetric in $a,b$, since interchanging them fixes
$s$ and negates $d$. The unique top-degree term gives~\eqref{eq:top}.
Finally there are $d\le L$ summands, so the error in~\eqref{eq:massgeneral}
becomes the error in~\eqref{eq:local}. This also proves uniformity across both
possible grid parities.
\end{proof}

\section{Why every lower asymptotic order cancels}

For decreasing offsets $a_1>\cdots>a_r\ge0$, define
\[
 \Delta(a)=\prod_{1\le i<j\le r}(a_i-a_j)>0.
\]
It has degree $R=\binom r2$. For integer distinct offsets, $\Delta(a)\ge1$.

\begin{proposition}\label{prop:localr}
Fix $r\ge2$ and $C>0$. There is a constant $\kappa_r$ such that, uniformly for
$0\le a_r<\cdots<a_1\le L\le C\log N$,
\begin{equation}\label{eq:localr}
 w_r(N-a_1,\ldots,N-a_r)
 =\kappa_r\Delta(a)N^{-R/2}
 +O_{r,C}\left((1+L)^{R+2}N^{-R/2-1}\right).
\end{equation}
The constant is the coefficient of $t^R\Delta(a)$ in the even or augmented
Pfaffian with ordinary entries $\operatorname{erf}(t(a_i-a_j)/2)$.
\end{proposition}
\begin{proof}
Let $m=\lfloor r/2\rfloor$ and set
\[
 K_r=\frac{R-m}{2}=
 \begin{cases}m(m-1),&r=2m,\\m^2,&r=2m+1.\end{cases}
\]
This is a nonnegative integer. Every matching in Lemma~\ref{lem:pf} has exactly
$m$ ordinary entries. In the odd case its one remaining entry is the constant
one associated with the dummy index.

Insert Lemma~\ref{lem:local} with $K=K_r$ into the Pfaffian and remove the common
factor $(\pi N)^{-m/2}$. Its coefficient at $N^{-k}$ is an alternating
polynomial in $a_1,\ldots,a_r$ of total degree at most $m+2k$. This follows
because a simultaneous permutation of ordinary rows and columns changes the
Pfaffian by the permutation sign. Every alternating polynomial over
$\mathbb Q$ is divisible by $\Delta(a)$: it vanishes when $a_i=a_j$, hence each
linear factor $a_i-a_j$ divides it, and these distinct prime factors have a
product dividing it. Thus all coefficients with $k<K_r$ vanish, since their
degree is below $R$. The coefficient at $k=K_r$ is $\lambda_r\Delta(a)$ for a
constant $\lambda_r$.

At that first possible order, only the top homogeneous parts~\eqref{eq:top}
can contribute. These are exactly the coefficients in
\[
 \operatorname{erf}\left(\frac{td}{2}\right)
 =\frac1{\sqrt\pi}\sum_{k\ge0}
       \frac{(-1)^k d^{2k+1}t^{2k+1}}{4^k k!(2k+1)}.
\]
Therefore $\kappa_r=\pi^{-m/2}\lambda_r$ is characterized as stated.

For clarity, the remainders are controlled before using any cancellation.
Each actual or truncated ordinary entry is $O_{r,C}((1+L)N^{-1/2})$.
Replacing one entry by the remainder from~\eqref{eq:local} gives, in a product
of $m$ entries,
\[
 O_{r,C}\left((1+L)^{2K_r+3+m-1}N^{-K_r-3/2-(m-1)/2}\right)
 =O_{r,C}\left((1+L)^{R+2}N^{-R/2-1}\right).
\]
The discarded terms of total index $k>K_r$ in the product of truncated
polynomials satisfy the same bound, since $(1+L)^2/N\to0$. There are only
finitely many pairings for fixed $r$. This proves~\eqref{eq:localr}.
\end{proof}

\section{Identification of the universal constant}

For completeness, set $\kappa_1=1$, as given by the augmented
Pfaffian when $r=1$.

\begin{lemma}\label{lem:kappa}
For each $r\ge1$,
\begin{equation}\label{eq:kappa}
 \kappa_r=
 \frac{\prod_{j=1}^r\Gamma(j/2)}
      {\pi^{r/2}\prod_{j=0}^{r-1}j!}>0.
\end{equation}
\end{lemma}
\begin{proof}
Fix any strictly decreasing real tuple $a_1>\cdots>a_r$, and let the independent
variables $Y_i$ be normal with mean $-ta_i$ and variance one. Their pairwise
sign expectations are $\operatorname{erf}(t(a_i-a_j)/2)$. Consequently the
Gaussian version of Lemma~\ref{lem:pf} identifies its Pfaffian as
\[
 F_a(t)=\E\sgn\prod_{i<j}(Y_j-Y_i).
\]
Let $\phi(x)=(2\pi)^{-1/2}e^{-x^2/2}$. Splitting $\R^r$ into its ordering
chambers gives
\begin{align}
 F_a(t)
 &=\int_{x_1<\cdots<x_r}\det[\phi(x_j+ta_i)]_{i,j=1}^r\,dx\notag\\
 &=e^{-t^2\sum_i a_i^2/2}
   \int_{x_1<\cdots<x_r}\det[e^{-ta_i x_j}]_{i,j=1}^r
           \prod_j\phi(x_j)\,dx.\label{eq:gaussian}
\end{align}
In the determinant of exponentials the first nonzero Taylor coefficient is
\[
 [t^R]\det[e^{-ta_i x_j}]
 =\frac{\Delta(a)\prod_{i<j}(x_j-x_i)}{\prod_{j=0}^{r-1}j!}.
\]
To verify this, expand each exponential in monomials. Distinct powers are
required in a nonzero determinant; their smallest possible sum is
$0+1+\cdots+(r-1)=R$. The corresponding two determinants are the Vandermonde
determinants in $-a_i$ and in $x_j$; the first is $\Delta(a)$.

Differentiation through the integrals to this finite order is valid. For $|t|\le1$
and fixed $a$, every derivative through order $R+1$ of the Gaussian densities is
bounded by a polynomial in $|x|$ times $C_a e^{-\sum_j x_j^2/4}$, which is
integrable. Thus if $Z_1,\ldots,Z_r$ are independent standard normals,
\begin{equation}\label{eq:kappaintegral}
 [t^R]F_a(t)=\frac{\Delta(a)}{r!\prod_{j=0}^{r-1}j!}
              \E\left|\prod_{i<j}(Z_j-Z_i)\right|.
\end{equation}
The factor $1/r!$ is the passage from one ordering chamber to all chambers.

The Gaussian Selberg (Mehta) integral at parameter $\gamma=1/2$ is
\begin{equation}\label{eq:Mehta}
 \E\left|\prod_{i<j}(Z_j-Z_i)\right|
 =\prod_{j=1}^r\frac{\Gamma(1+j/2)}{\Gamma(3/2)}.
\end{equation}
This is the normalization in~\cite[eqs.~(1.5)--(1.8)]{fw}. Its hypotheses hold:
$r$ is a positive integer, $\gamma=1/2$ is in the positive-parameter domain,
the integrand has no nonintegrable collision singularities, and its tails are
a fixed polynomial times a Gaussian. Combining~\eqref{eq:kappaintegral} and
\eqref{eq:Mehta}, then using $\Gamma(1+j/2)=(j/2)\Gamma(j/2)$ and
$\Gamma(3/2)=\sqrt\pi/2$, proves~\eqref{eq:kappa}. All gamma arguments are
positive, so $\kappa_r>0$.
\end{proof}

\section{The weighted discrete Vandermonde sum}

\begin{lemma}\label{lem:Z}
For a positive integer $r$ and $0<q<1$, define
\[
 Z_r(q)=\sum_{0\le b_1<\cdots<b_r}
          q^{\sum_i b_i}\prod_{i<j}(b_j-b_i)^2.
\]
Then
\begin{equation}\label{eq:Z}
 Z_r(q)=\frac{q^{\binom r2}\prod_{j=0}^{r-1}(j!)^2}{(1-q)^{r^2}}.
\end{equation}
\end{lemma}
\begin{proof}
Let $\mu_k(q)=\sum_{b\ge0}b^kq^b$, where $0^0=1$. All these moments converge.
Cauchy--Binet applied first to a finite weighted Vandermonde matrix, then taking
its number of columns to infinity, gives
\[
 Z_r(q)=\det[\mu_{i+j}(q)]_{0\le i,j<r}.
\]
The passage to infinity is legitimate because each moment converges and a finite
determinant is a polynomial in its entries.

For an elementary evaluation of this determinant, define polynomials $M_j(x)$ by
\[
 \sum_{j\ge0}M_j(x)z^j
 =\frac{(1-z/q)^x}{(1-z)^{x+1}},\qquad x\in\Z_{\ge0}.
\]
The degree of $M_j$ is $j$ and its leading coefficient is
$(1-1/q)^j/j!$. Summing a geometric series gives
\begin{align*}
 &\sum_{x\ge0}q^x
 \frac{(1-z/q)^x}{(1-z)^{x+1}}
 \frac{(1-w/q)^x}{(1-w)^{x+1}}\\
 &\qquad=\frac1{(1-z)(1-w)-q(1-z/q)(1-w/q)}
 =\frac1{(1-q)(1-zw/q)}.
\end{align*}
For $z,w$ in a sufficiently small neighborhood of zero the geometric series
converges absolutely and uniformly on smaller neighborhoods, since its ratio
is $q<1$ at zero. Coefficient comparison is therefore valid and yields
\[
 \sum_{x\ge0}q^xM_j(x)M_k(x)=\frac{q^{-j}}{1-q}\,\delta_{jk}.
\]
After making the polynomials monic, their squared norms are
\[
 h_j=\frac{(j!)^2q^j}{(1-q)^{2j+1}}.
\]
The change from monomials to monic polynomials is unit triangular, so the moment
determinant is $\prod_{j=0}^{r-1}h_j$. This is~\eqref{eq:Z}.
\end{proof}

\section{Tail control and completion of the theorem}

\begin{proof}[Proof of Theorem~\ref{thm:main}]
The $r=0$ case is immediate from the distance constraint. For $r=1$, the
sum-of-squares identity gives
\[
 H_1(n,n)=\sum_{u=0}^{n-2}4^u=\frac{4^{n-1}-1}{3},
\]
with the empty sum at $n=1$. This has the stated constant $1/12$.

Fix now $r\ge2$. Replacing row indices by decreasing offsets in
\eqref{eq:squares} gives
\begin{equation}\label{eq:normalH}
 \frac{N^RH_r(n,n)}{4^{rN}}
 =\sum_{N\ge a_1>\cdots>a_r\ge0}
       4^{-\sum_i a_i}N^Rw_r(N-a_1,\ldots,N-a_r)^2.
\end{equation}
Choose
\[
 L=\left\lceil(R+2)\log_4 N\right\rceil.
\]
The window $L=O_r(\log N)$ is covered by Proposition~\ref{prop:localr} for a
fixed $C$ depending only on $r$. Outside it, the global bound $|w_r|\le1$ gives
\begin{align}
 &N^R\sum_{\max_i a_i>L}4^{-\sum_i a_i}\notag\\
 &\quad\le rN^R\left(\sum_{a>L}4^{-a}\right)
                     \left(\sum_{a\ge0}4^{-a}\right)^{r-1}
 =\frac r3\left(\frac43\right)^{r-1}N^R4^{-L}
 =O_r(N^{-2}).\label{eq:tail}
\end{align}
Here the sum was enlarged to all nonnegative integer tuples, so it also bounds
all allowable decreasing tuples.

Inside the window, $\Delta(a)\ge1$ and $\kappa_r>0$. The relative error in
\eqref{eq:localr} is consequently
$O_r((1+L)^{R+2}/N)=o(1)$ uniformly. Squaring gives
\[
 N^Rw_r(N-a_1,\ldots,N-a_r)^2
 =\kappa_r^2\Delta(a)^2(1+o(1))
\]
uniformly there. The weighted main sum converges to $Z_r(1/4)$, which is finite
by Lemma~\ref{lem:Z}. The uniform relative error therefore contributes $o(1)$
to the sum. Combining this with~\eqref{eq:tail} yields
\begin{equation}\label{eq:Nlimit}
 H_r(n,n)\sim4^{rN}N^{-R}\kappa_r^2 Z_r(1/4).
\end{equation}
No exchange of a merely pointwise limit with an uncontrolled infinite sum is
used.

Finally $N=n-2$, and
\begin{align*}
 c_r
 &=4^{-2r}\kappa_r^2 Z_r(1/4)\\
 &=4^{-2r}\,
 \frac{\prod_{j=1}^r\Gamma(j/2)^2}{\pi^r\prod_{j=0}^{r-1}(j!)^2}
 \frac{4^{r^2-R}\prod_{j=0}^{r-1}(j!)^2}{3^{r^2}}\\
 &=\frac{2^{r(r-3)}}{3^{r^2}\pi^r}\prod_{j=1}^r\Gamma(j/2)^2.
\end{align*}
We used $N/n\to1$ and $r^2-R-2r=r(r-3)/2$. This proves~\eqref{eq:main}.
For $r=4$, the gamma product before squaring is
$\Gamma(1/2)\Gamma(1)\Gamma(3/2)\Gamma(2)=\pi/2$, giving
$c_4=4/(3^{16}\pi^2)$ exactly.
\end{proof}

\section{A direct four-row calculation}

This section gives an independent, elementary extraction of the required
four-row coefficient. In particular, the $r=4$ conclusion does not depend on
the general Gaussian integral evaluation in Lemma~\ref{lem:kappa}.

Put $d=a-b$ and $s=a+b$. Lemma~\ref{lem:local} through $K=2$ reads
\begin{align}
 Q(N-a,N-b)
 &=\frac d{\sqrt{\pi N}}
   \left(1+\frac{A_1(s,d)}{N}+\frac{A_2(s,d)}{N^2}\right)
 +O\left((1+L)^7N^{-7/2}\right),\label{eq:Qexplicit}\\
 A_1(s,d)&=\frac s4-\frac{d^2}{12}-\frac7{24},\notag\\
 A_2(s,d)&=\frac{d^4}{160}-\frac{sd^2}{16}
       +\frac{3s^2}{32}+\frac{3d^2}{32}-\frac{7s}{32}+\frac{61}{640}.\notag
\end{align}
For example the mass at $j=m/2+t$ has first two coefficients
\begin{align*}
 T_1(s,t)&=s/4-t^2-1/8,\\
 T_2(s,t)&=3s^2/32-3st^2/4-3s/32+t^4/2+5t^2/8+1/128.
\end{align*}
Summing these and using
\[
 \sum t^2=\frac{d(d^2+2)}{12},\qquad
 \sum t^4=\frac{d(3d^4+20d^2-8)}{240}
\]
gives~\eqref{eq:Qexplicit}.

Insert it into~\eqref{eq:pf4}. After the common factor $(\pi N)^{-1}$,
the coefficients of $N^0$ and $N^{-1}$ are alternating polynomials of degrees
at most two and four. Both vanish. The coefficient of $N^{-2}$ is an
alternating polynomial of degree at most six, hence a constant times
$\Delta(a_1,a_2,a_3,a_4)$. Only the leading terms
\[
 g(d)=d-\frac{d^3}{12}\varepsilon+\frac{d^5}{160}\varepsilon^2
\]
are needed to determine that constant. At $(a_1,a_2,a_3,a_4)=(3,2,1,0)$,
\[
 [\varepsilon^2]\bigl(g(1)^2-g(2)^2+g(3)g(1)\bigr)=\frac12,
 \qquad \Delta(3,2,1,0)=12.
\]
Thus the constant is $1/24$, and
\begin{equation}\label{eq:local4explicit}
 w_4(N-a_1,\ldots,N-a_4)
 =\frac{\Delta(a)}{24\pi N^3}
  +O\left((1+L)^8N^{-4}\right).
\end{equation}
The remaining constant is
\[
 Z_4(q)=\frac{144q^6}{(1-q)^{16}},\qquad
 Z_4(1/4)=\frac{144\cdot4^{10}}{3^{16}}.
\]
Equations~\eqref{eq:Nlimit} and~\eqref{eq:local4explicit} give
\[
 H_4(n,n)\sim\frac{256^N}{576\pi^2N^6}
                   \frac{144\cdot4^{10}}{3^{16}}.
\]
Replacing $N$ by $n-2$ contributes the essential factor
$256^{-2}=4^{-8}$. The coefficient becomes
\[
 \frac{144\cdot4^{10}}{576\cdot4^8\cdot3^{16}\pi^2}
 =\frac4{3^{16}\pi^2},
\]
which independently verifies the exact target constant.

\section{Constants, source coverage and limits of scope}

Theorem~\ref{thm:main} holds for each fixed integer $r\ge1$, with
$H_0(n,n)=1$ separately. It makes no uniform assertion for $r$ growing with
$n$ and no assertion about joint rectangular limits. The exponent and
constants agree with all rows through $r=5$ in~\cite[p.~12]{dk}:
\[
\begin{array}{c|ccccc}
r&1&2&3&4&5\\[2pt]\hline
c_r&\dfrac1{12}&\dfrac1{2^2 3^4\pi}&\dfrac1{2^2 3^9\pi}
&\dfrac{2^2}{3^{16}\pi^2}&\dfrac{2^4}{3^{23}\pi^2}
\end{array}
\]
The formula also gives
\[
 c_6=\frac{2^{14}}{3^{34}\pi^3},\qquad
 c_7=\frac{2^{18}\,5^2}{3^{45}\pi^3}.
\]
Thus the source's further suggestion that all constants use only powers of
$2,3,\pi$ cannot hold with integer exponents for every $r$. The factor $5^2$
at $r=7$ cannot be absorbed into those powers, by the transcendence of $\pi$
and unique prime factorization. This concerns a separate informal general
pattern: it does not contradict the displayed rows of the source table or
the fixed-parameter asymptotic scale. Allowing arbitrary real exponents
would remove the intended prime-factor restriction.

\paragraph{Related rectangular work.}
Dougherty-Bliss and Spahn's published four-page ISSAC poster abstract
\emph{Rectangular Hardinian Arrays}~\cite{ds}, in ACM Communications in
Computer Algebra, is available and has been read for this comparison.
For fixed $r$ and fixed width $k$, its theorem makes $H_r(n,k)$ a polynomial
of degree $r$ in $n$ for sufficiently large $n$, using a finite transfer
matrix. Its hypotheses keep $k$ fixed; they do not permit substituting
$k=n$ or give a uniform threshold and asymptotic estimate in growing $k$.
That theorem therefore does not supply the square-array result proved here.
This is a comparison of specified sources, not certification of worldwide
publication priority.

\paragraph{Proof dependencies and finite evidence.}
The original complementary-minor count is an external combinatorial input
from~\cite{dk}. The square-sum and sign-expectation reductions from the
public note~\cite{prior} are rederived here with all conventions fixed.
Jacobi's identity, Cauchy--Binet, Stirling's formula with its stated remainder,
Faulhaber's identity and the Mehta integral are used explicitly at their
respective steps. In particular, the all-$r$ constant uses the standard
Gaussian integral~\cite{fw}, whereas the four-row calculation extracts its
coefficient directly.

No finite computation is needed to conclude the theorem. Earlier records
contain both an author-reported exact-data run and separate later finite
checks; the distribution's reproducibility note distinguishes these records
from this manuscript preparation, during which no mathematical program was
run. File hashes identify versions and do not themselves certify execution.
None of these finite checks replaces the arbitrary-order expansion,
Vandermonde cancellation, integral normalization or full tail estimate.

\paragraph{Version and verification status.}
This is the fixed-parameter extension volume of the public $r=2,3$
note~\cite{prior}; the earlier manuscript, code and data remain unchanged
historical material. Prepared with AI assistance. No independent human
peer-review, proof-assistant certification or worldwide-priority claim is
made.

\begin{thebibliography}{9}
\raggedright
\bibitem{dk}
R.~Dougherty-Bliss and M.~Kauers,
\emph{Hardinian Arrays}, Electronic Journal of Combinatorics
\textbf{31}(2) (2024), P2.9, 13 pages.
\href{https://doi.org/10.37236/12358}{doi:10.37236/12358}.
Theorem~6, p.~9; asymptotic table, p.~12.

\bibitem{prior}
Aurantium, \emph{Two asymptotic formulas for square Hardinian arrays},
public research note, version~0.1.0 (2026).
\href{https://github.com/Aurantium-47/hardinian-array-asymptotics/releases/tag/v0.1.0}{Public version and source archive}.
Fixed source commit
\href{https://github.com/Aurantium-47/hardinian-array-asymptotics/blob/8bf9f116a7d49c6b49ecfec9d30d14a2afc4502e/main.tex}{8bf9f116a7d49c6b49ecfec9d30d14a2afc4502e}.
Section~2 derives the general square-sum and probability identities; the
asymptotic conclusions of that volume cover $r=2,3$.

\bibitem{fw}
P.~J.~Forrester and S.~O.~Warnaar,
\emph{The importance of the Selberg integral}, Bulletin of the American
Mathematical Society (N.S.) \textbf{45} (2008), 489--534.
\href{https://arxiv.org/abs/0710.3981}{arXiv:0710.3981}.
The Gaussian integral normalization and evaluation are (1.5)--(1.8).

\bibitem{dlmf}
NIST Digital Library of Mathematical Functions,
\href{https://dlmf.nist.gov/5.11}{\S5.11(i)--(ii)},
logarithmic Stirling expansion and positive-real-axis remainder bounds.

\bibitem{ds}
R.~Dougherty-Bliss and G.~Spahn,
\emph{Rectangular Hardinian Arrays}, ACM Communications in Computer Algebra
\textbf{58}(3), issue~229 (September 2024), 81--84,
ISSAC 2024 poster abstracts.
\href{https://doi.org/10.1145/3717582.3717589}{doi:10.1145/3717582.3717589}.
\end{thebibliography}
\end{document}
